\section{位置は二種類ある}
\label{sec:position}

地位は\dstrec{}である。多体の層力と、関係としての位置は、別の型に入れる。

\subsection{チャート上の点}

\texttt{LayerForce} の位置は、実内積空間 \(E\) の点である。三次元の計算では \(E=\mathbb{R}^3\) で足りる。粒子の識別子を \(\iota\) とすると、配置は
\begin{verbatim}
pos:  map iota -> Vec3
\end{verbatim}
である。分離は \(\lVert\mathrm{pos}(i)-\mathrm{pos}(j)\rVert\) であり、このノルムが層の引数 \(r\) になる。この \(E\) は粒子の内在角が住む代数ではない。

\subsection{三球上の関係}

\texttt{RelationalPosition.S3} は単位四元数である。
\begin{verbatim}
struct S3:
    val: quaternion
    # normSq(val) = 1
\end{verbatim}
\(b\) の \(a\) に対する位置は相対ローター
\[
\rho(a,b)=a^{-1}b
\]
であり、識別子は \texttt{relative} である。\(\rho(a,b)=1\) は \(a=b\) と同値である。左からの同時移動は \(\rho\) を変えない。

この \(\rho\) は、粒子内部の不整合ローター \(\Omega\) ではない。チャート座標の位置演算子でもない。\(\rho\) を \(\mathbb{R}^3\) の差と見なして、層力の \(r\) に渡さない。\(S^3\) 上の測地線を、ニュートン軌道の代わりに積分しない。

\subsection{状態を分ける}

最小の実行時状態は次の三層に割る。

\begin{verbatim}
struct ChartBody:
    id: iota
    pos: Vec3
    mass_input: real    # factor in G * m_i * m_j
    charge_input: real  # factor in k * q_i * q_j

struct WrittenSeed:
    axis: int           # 0, 1, or 2
    sigma0: real
    sigmadot0: real
    delta0: real
    nu: real            # constant delta-dot
    stiffness: real     # input m

struct RelationalSample:
    element: S3         # not read by layer_force
\end{verbatim}

\texttt{ChartBody} は第\ref{sec:force}章だけが読む。\texttt{WrittenSeed} は第\ref{sec:written}章だけが読む。\texttt{RelationalSample} は関係の検査に残し、時間ループの位置更新へ接続しない。一つの構造体に三層の場を並べることはできる。そのとき更新関数の引数リストで、読む場を制限する。
